%%%%%%%%%%  scan idx 112  |  folio 101  %%%%%%%%%%

\medskip
\noindent\textbf{Observación 4.6(a).}
Sea $X_0$ afín. Existe una deformación formal versal de $X_0$ si y
solo si el soporte de $\underline D^1_{X_0}$ es un conjunto finito.
Esto no implica, sin embargo, que todos los puntos no lisos de $X_0$
sean aislados: existen esquemas afines $X_0$ sobre $k$ que poseen un
punto no liso no aislado y para los cuales
$\operatorname{Supp}\underline D^1_{X_0}$ es finito.

\medskip
\noindent\textbf{Observación 4.6(b).}
En lo sucesivo, una deformación de $X_0$ significará un objeto de
$\underline M_{X_0}$. Cuando pueda haber confusión, tal objeto se
llamará una deformación infinitesimal de $X_0$.

\medskip
\noindent\textbf{Observación 4.6(c).}
Sea $X_0$ un esquema sobre $k$. Hemos fijado un anillo local
noetheriano completo $\Lambda$ con cuerpo residual $k$ y considerado
las deformaciones de $X_0$ sobre $\underline C_\Lambda$.

%%%%%%%%%%  scan idx 113  |  folio 102  %%%%%%%%%%

En la práctica se toma habitualmente $\Lambda=k$, o bien
$\Lambda=W(k)$ cuando $k$ es un cuerpo perfecto de característica
positiva. La inclusión
$\underline C_k\subset\underline C_\Lambda$ es plenamente fiel, y
$\underline M_{X_0}|_{\underline C_k}$ es simplemente la categoría
de deformaciones de $X_0$ sobre los objetos de $\underline C_k$.
Por consiguiente, si $(R,\mathfrak X)$ es una deformación formal
versal de $X_0$ sobre $\underline C_\Lambda$, entonces
\[
 (k\otimes_\Lambda R,\;k\otimes_\Lambda\mathfrak X)
\]
es una deformación formal versal de $X_0$ sobre
$\underline C_k$. Se sigue, por ejemplo, que
\[
 \dim_{\underline C_\Lambda}\underline M_{X_0}
 =\dim(k\otimes_\Lambda R)
\]
no depende de la elección de $\Lambda$, sino únicamente de $X_0$.

El teorema precedente admite muchas generalizaciones. Por ejemplo,
consideremos un par $Y_0\longrightarrow X_0$, donde $X_0$ es un
esquema de tipo finito sobre $k$ y $Y_0$ es un subesquema cerrado fijo
de $X_0$. Para
$R\in\operatorname{Ob}(\underline C_\Lambda)$, una deformación de
$(X_0,Y_0)$ sobre $R$ es un diagrama conmutativo
\[
\begin{tikzcd}[row sep=large,column sep=huge]
Y_0 \arrow[rrr] \arrow[dr,hook] \arrow[ddr]
&&& Y \arrow[dl,hook] \arrow[ddl,"\mathrm{flat}"] \\
& X_0 \arrow[r] \arrow[d]
& X \sisi{\arrow[d,no head,"\mathrm{flat}"']}{\arrow[d,"\mathrm{flat}"']} & \\
& \operatorname{Spec}(k) \arrow[r]
& \operatorname{Spec}(R) &
\end{tikzcd}
\]
tal que la aplicación inducida
\[
(X_0,Y_0)\longrightarrow
\left(
 X\times_{\operatorname{Spec}(R)}\operatorname{Spec}(k),
 Y\times_{\operatorname{Spec}(R)}\operatorname{Spec}(k)
\right)
\]
sea un isomorfismo. Definir de la manera evidente los morfismos
\[
(R';X',Y')\longrightarrow(R;X,Y)
\]
de deformaciones de $(X_0,Y_0)$ proporciona una categoría cofibrada
$\underline M_{X_0,Y_0}$ en grupoides sobre
$\underline C_\Lambda$. Esta construcción se utiliza con frecuencia
para rigidificar los automorfismos en problemas de módulos de
deformaciones. El argumento del Corolario 4.3, basado en el Lema 4.2,
muestra inmediatamente que $\underline M_{X_0,Y_0}$ es una categoría
cofibrada homogénea en grupoides sobre $\underline C_\Lambda$.
